% =====================================================================
% TEXTBOOK BIB — verbose source metadata for every paper in the project.
% One entry per paper: identification, abstract, keywords, contents,
% references. All content verbatim from the source documents.
% =====================================================================
\documentclass[11pt,a4paper]{article}

\newcommand{\papertag}{textbook\_bib}

% =====================================================================
% Shared preamble for paper notes
% Include with:  % =====================================================================
% Shared preamble for paper notes
% Include with:  % =====================================================================
% Shared preamble for paper notes
% Include with:  \input{../_common/preamble}
% =====================================================================

% --- Core packages ---
\usepackage[utf8]{inputenc}
\usepackage[T1]{fontenc}
\usepackage{lmodern}
\usepackage[margin=2.5cm]{geometry}
\usepackage{amsmath,amssymb,amsthm,mathtools}
\usepackage{bm}
\usepackage{booktabs}
\usepackage{array}
\usepackage{enumitem}
\usepackage{hyperref}
\usepackage{xcolor}
\usepackage{listings}
\usepackage{fancyhdr}
\usepackage{titlesec}

% --- Hyperref styling ---
\hypersetup{
  colorlinks=true,
  linkcolor=blue!50!black,
  urlcolor=blue!50!black,
  citecolor=blue!50!black
}

% --- Theorem-like environments ---
\theoremstyle{definition}
\newtheorem{defn}{Definition}
\newtheorem{remark}{Remark}
\newtheorem{example}{Example}

\theoremstyle{plain}
\newtheorem{prop}{Proposition}
\newtheorem{lemma}{Lemma}

% --- Coloured boxes for the structured sections ---
% Validation block, commentary, TODOs use coloured frames so the eye
% can find them when flipping through the note.
\usepackage[most]{tcolorbox}

\newtcolorbox{commentary}{
  colback=yellow!5,
  colframe=yellow!50!black,
  title=Commentary,
  fonttitle=\bfseries,
  breakable
}

\newtcolorbox{validation}{
  colback=green!3,
  colframe=green!50!black,
  title=Validation,
  fonttitle=\bfseries,
  breakable
}

\newtcolorbox{todobox}{
  colback=red!3,
  colframe=red!50!black,
  title=TODO / Open questions,
  fonttitle=\bfseries,
  breakable
}

% --- Code listing style for Python snippets in validation blocks ---
\lstdefinestyle{py}{
  language=Python,
  basicstyle=\ttfamily\footnotesize,
  keywordstyle=\color{blue!60!black},
  commentstyle=\color{gray},
  stringstyle=\color{red!60!black},
  showstringspaces=false,
  breaklines=true,
  frame=single,
  framesep=4pt,
  rulecolor=\color{gray!30}
}

% --- Page header: shows paper tag so a printed note is identifiable ---
\pagestyle{fancy}
\fancyhf{}
\fancyhead[L]{\small\textit{\papertag}}
\fancyhead[R]{\small\thepage}
\renewcommand{\headrulewidth}{0.4pt}

% --- Section spacing: tighter than article default ---
\titlespacing*{\section}{0pt}{1.5ex plus 0.5ex minus 0.2ex}{1ex plus 0.2ex}
\titlespacing*{\subsection}{0pt}{1.2ex plus 0.3ex minus 0.2ex}{0.6ex plus 0.2ex}

% =====================================================================

% --- Core packages ---
\usepackage[utf8]{inputenc}
\usepackage[T1]{fontenc}
\usepackage{lmodern}
\usepackage[margin=2.5cm]{geometry}
\usepackage{amsmath,amssymb,amsthm,mathtools}
\usepackage{bm}
\usepackage{booktabs}
\usepackage{array}
\usepackage{enumitem}
\usepackage{hyperref}
\usepackage{xcolor}
\usepackage{listings}
\usepackage{fancyhdr}
\usepackage{titlesec}

% --- Hyperref styling ---
\hypersetup{
  colorlinks=true,
  linkcolor=blue!50!black,
  urlcolor=blue!50!black,
  citecolor=blue!50!black
}

% --- Theorem-like environments ---
\theoremstyle{definition}
\newtheorem{defn}{Definition}
\newtheorem{remark}{Remark}
\newtheorem{example}{Example}

\theoremstyle{plain}
\newtheorem{prop}{Proposition}
\newtheorem{lemma}{Lemma}

% --- Coloured boxes for the structured sections ---
% Validation block, commentary, TODOs use coloured frames so the eye
% can find them when flipping through the note.
\usepackage[most]{tcolorbox}

\newtcolorbox{commentary}{
  colback=yellow!5,
  colframe=yellow!50!black,
  title=Commentary,
  fonttitle=\bfseries,
  breakable
}

\newtcolorbox{validation}{
  colback=green!3,
  colframe=green!50!black,
  title=Validation,
  fonttitle=\bfseries,
  breakable
}

\newtcolorbox{todobox}{
  colback=red!3,
  colframe=red!50!black,
  title=TODO / Open questions,
  fonttitle=\bfseries,
  breakable
}

% --- Code listing style for Python snippets in validation blocks ---
\lstdefinestyle{py}{
  language=Python,
  basicstyle=\ttfamily\footnotesize,
  keywordstyle=\color{blue!60!black},
  commentstyle=\color{gray},
  stringstyle=\color{red!60!black},
  showstringspaces=false,
  breaklines=true,
  frame=single,
  framesep=4pt,
  rulecolor=\color{gray!30}
}

% --- Page header: shows paper tag so a printed note is identifiable ---
\pagestyle{fancy}
\fancyhf{}
\fancyhead[L]{\small\textit{\papertag}}
\fancyhead[R]{\small\thepage}
\renewcommand{\headrulewidth}{0.4pt}

% --- Section spacing: tighter than article default ---
\titlespacing*{\section}{0pt}{1.5ex plus 0.5ex minus 0.2ex}{1ex plus 0.2ex}
\titlespacing*{\subsection}{0pt}{1.2ex plus 0.3ex minus 0.2ex}{0.6ex plus 0.2ex}

% =====================================================================

% --- Core packages ---
\usepackage[utf8]{inputenc}
\usepackage[T1]{fontenc}
\usepackage{lmodern}
\usepackage[margin=2.5cm]{geometry}
\usepackage{amsmath,amssymb,amsthm,mathtools}
\usepackage{bm}
\usepackage{booktabs}
\usepackage{array}
\usepackage{enumitem}
\usepackage{hyperref}
\usepackage{xcolor}
\usepackage{listings}
\usepackage{fancyhdr}
\usepackage{titlesec}

% --- Hyperref styling ---
\hypersetup{
  colorlinks=true,
  linkcolor=blue!50!black,
  urlcolor=blue!50!black,
  citecolor=blue!50!black
}

% --- Theorem-like environments ---
\theoremstyle{definition}
\newtheorem{defn}{Definition}
\newtheorem{remark}{Remark}
\newtheorem{example}{Example}

\theoremstyle{plain}
\newtheorem{prop}{Proposition}
\newtheorem{lemma}{Lemma}

% --- Coloured boxes for the structured sections ---
% Validation block, commentary, TODOs use coloured frames so the eye
% can find them when flipping through the note.
\usepackage[most]{tcolorbox}

\newtcolorbox{commentary}{
  colback=yellow!5,
  colframe=yellow!50!black,
  title=Commentary,
  fonttitle=\bfseries,
  breakable
}

\newtcolorbox{validation}{
  colback=green!3,
  colframe=green!50!black,
  title=Validation,
  fonttitle=\bfseries,
  breakable
}

\newtcolorbox{todobox}{
  colback=red!3,
  colframe=red!50!black,
  title=TODO / Open questions,
  fonttitle=\bfseries,
  breakable
}

% --- Code listing style for Python snippets in validation blocks ---
\lstdefinestyle{py}{
  language=Python,
  basicstyle=\ttfamily\footnotesize,
  keywordstyle=\color{blue!60!black},
  commentstyle=\color{gray},
  stringstyle=\color{red!60!black},
  showstringspaces=false,
  breaklines=true,
  frame=single,
  framesep=4pt,
  rulecolor=\color{gray!30}
}

% --- Page header: shows paper tag so a printed note is identifiable ---
\pagestyle{fancy}
\fancyhf{}
\fancyhead[L]{\small\textit{\papertag}}
\fancyhead[R]{\small\thepage}
\renewcommand{\headrulewidth}{0.4pt}

% --- Section spacing: tighter than article default ---
\titlespacing*{\section}{0pt}{1.5ex plus 0.5ex minus 0.2ex}{1ex plus 0.2ex}
\titlespacing*{\subsection}{0pt}{1.2ex plus 0.3ex minus 0.2ex}{0.6ex plus 0.2ex}

% =====================================================================
% House notation: shared symbols used across all paper notes.
% Each note imports this and may shadow individual macros if a paper
% uses a clashing convention — but the *default* meaning is fixed here.
%
% Include with:  % =====================================================================
% House notation: shared symbols used across all paper notes.
% Each note imports this and may shadow individual macros if a paper
% uses a clashing convention — but the *default* meaning is fixed here.
%
% Include with:  % =====================================================================
% House notation: shared symbols used across all paper notes.
% Each note imports this and may shadow individual macros if a paper
% uses a clashing convention — but the *default* meaning is fixed here.
%
% Include with:  \input{../_common/notation}
% =====================================================================

% --- Measures and expectations ---
\newcommand{\Q}{\mathbb{Q}}              % risk-neutral measure
\newcommand{\Qt}{\mathbb{Q}^{T}}         % T-forward measure
\newcommand{\E}{\mathbb{E}}              % expectation
\newcommand{\Et}[1]{\E^{#1}}             % expectation under measure #1
\newcommand{\Ftime}[1]{\mathcal{F}_{#1}} % filtration at #1

% --- Discount and forward objects ---
\newcommand{\Disc}[2]{D_{#1,#2}}         % discount factor D_{t,T}
\newcommand{\Bond}[2]{P_{#1}(#2)}        % bond price P_t(y)
\newcommand{\Ann}[1]{A_{#1}}             % annuity / level
\newcommand{\Numer}[1]{N_{#1}}           % numeraire
\newcommand{\Fwd}[2]{F_{#1,#2}}          % forward price F_{t,T}
\newcommand{\Fwdpx}[2]{\mathrm{Fwd}_{#1}(#2)} % verbose forward-price F_t(T)

% --- Rates ---
\newcommand{\IRR}{\mathrm{IRR}}          % internal rate of return
\newcommand{\Yld}{y}                     % bond yield variable
\newcommand{\Strike}{K}                  % strike (price-space)
\newcommand{\Lock}{L}                    % locked yield / rate
\newcommand{\Repo}{r^{\mathrm{repo}}}    % repo rate
\newcommand{\Fund}{r^{\mathrm{fun}}}     % unsecured funding rate
\newcommand{\OIS}{r^{\mathrm{ois}}}      % OIS rate
\newcommand{\Coll}{r^{\mathrm{c}}}       % collateral rate
\newcommand{\Rcmt}{R^{\mathrm{cmt}}}     % constant-maturity-treasury rate
\newcommand{\Rswp}{R^{\mathrm{swp}}}     % swap rate (used for risky variant)
\newcommand{\Rec}{\mathrm{Rec}}          % recovery rate
\newcommand{\NPV}{\mathrm{NPV}}          % present-value functional
\newcommand{\PVone}{\mathrm{PV01}}       % risky annuity
\newcommand{\Loss}{\mathrm{Loss}}        % default-leg integral
\newcommand{\Sasw}{S^{\mathrm{ASW}}}     % asset-swap spread
\newcommand{\Scds}{S^{\mathrm{CDS}}}     % CDS spread
\newcommand{\Basis}{\mathrm{Basis}}      % CDS-bond basis

% --- Sensitivities ---
\newcommand{\Diff}[2]{D_{#1}^{#2}}       % Euler's notation: D_x^k[f]
\newcommand{\Riskf}{\mathrm{RiskFactor}} % bond risk factor (= -DV01*1e4)
\newcommand{\DVone}{\mathrm{DV01}}

% --- Indicator / misc ---
\newcommand{\Ind}{\mathbf{1}}
\newcommand{\given}{\,|\,}
\newcommand{\dd}{\,\mathrm{d}}

% --- Greeks-style ---
\newcommand{\Gam}{\Gamma}
\newcommand{\Vega}{\mathcal{V}}

% =====================================================================

% --- Measures and expectations ---
\newcommand{\Q}{\mathbb{Q}}              % risk-neutral measure
\newcommand{\Qt}{\mathbb{Q}^{T}}         % T-forward measure
\newcommand{\E}{\mathbb{E}}              % expectation
\newcommand{\Et}[1]{\E^{#1}}             % expectation under measure #1
\newcommand{\Ftime}[1]{\mathcal{F}_{#1}} % filtration at #1

% --- Discount and forward objects ---
\newcommand{\Disc}[2]{D_{#1,#2}}         % discount factor D_{t,T}
\newcommand{\Bond}[2]{P_{#1}(#2)}        % bond price P_t(y)
\newcommand{\Ann}[1]{A_{#1}}             % annuity / level
\newcommand{\Numer}[1]{N_{#1}}           % numeraire
\newcommand{\Fwd}[2]{F_{#1,#2}}          % forward price F_{t,T}
\newcommand{\Fwdpx}[2]{\mathrm{Fwd}_{#1}(#2)} % verbose forward-price F_t(T)

% --- Rates ---
\newcommand{\IRR}{\mathrm{IRR}}          % internal rate of return
\newcommand{\Yld}{y}                     % bond yield variable
\newcommand{\Strike}{K}                  % strike (price-space)
\newcommand{\Lock}{L}                    % locked yield / rate
\newcommand{\Repo}{r^{\mathrm{repo}}}    % repo rate
\newcommand{\Fund}{r^{\mathrm{fun}}}     % unsecured funding rate
\newcommand{\OIS}{r^{\mathrm{ois}}}      % OIS rate
\newcommand{\Coll}{r^{\mathrm{c}}}       % collateral rate
\newcommand{\Rcmt}{R^{\mathrm{cmt}}}     % constant-maturity-treasury rate
\newcommand{\Rswp}{R^{\mathrm{swp}}}     % swap rate (used for risky variant)
\newcommand{\Rec}{\mathrm{Rec}}          % recovery rate
\newcommand{\NPV}{\mathrm{NPV}}          % present-value functional
\newcommand{\PVone}{\mathrm{PV01}}       % risky annuity
\newcommand{\Loss}{\mathrm{Loss}}        % default-leg integral
\newcommand{\Sasw}{S^{\mathrm{ASW}}}     % asset-swap spread
\newcommand{\Scds}{S^{\mathrm{CDS}}}     % CDS spread
\newcommand{\Basis}{\mathrm{Basis}}      % CDS-bond basis

% --- Sensitivities ---
\newcommand{\Diff}[2]{D_{#1}^{#2}}       % Euler's notation: D_x^k[f]
\newcommand{\Riskf}{\mathrm{RiskFactor}} % bond risk factor (= -DV01*1e4)
\newcommand{\DVone}{\mathrm{DV01}}

% --- Indicator / misc ---
\newcommand{\Ind}{\mathbf{1}}
\newcommand{\given}{\,|\,}
\newcommand{\dd}{\,\mathrm{d}}

% --- Greeks-style ---
\newcommand{\Gam}{\Gamma}
\newcommand{\Vega}{\mathcal{V}}

% =====================================================================

% --- Measures and expectations ---
\newcommand{\Q}{\mathbb{Q}}              % risk-neutral measure
\newcommand{\Qt}{\mathbb{Q}^{T}}         % T-forward measure
\newcommand{\E}{\mathbb{E}}              % expectation
\newcommand{\Et}[1]{\E^{#1}}             % expectation under measure #1
\newcommand{\Ftime}[1]{\mathcal{F}_{#1}} % filtration at #1

% --- Discount and forward objects ---
\newcommand{\Disc}[2]{D_{#1,#2}}         % discount factor D_{t,T}
\newcommand{\Bond}[2]{P_{#1}(#2)}        % bond price P_t(y)
\newcommand{\Ann}[1]{A_{#1}}             % annuity / level
\newcommand{\Numer}[1]{N_{#1}}           % numeraire
\newcommand{\Fwd}[2]{F_{#1,#2}}          % forward price F_{t,T}
\newcommand{\Fwdpx}[2]{\mathrm{Fwd}_{#1}(#2)} % verbose forward-price F_t(T)

% --- Rates ---
\newcommand{\IRR}{\mathrm{IRR}}          % internal rate of return
\newcommand{\Yld}{y}                     % bond yield variable
\newcommand{\Strike}{K}                  % strike (price-space)
\newcommand{\Lock}{L}                    % locked yield / rate
\newcommand{\Repo}{r^{\mathrm{repo}}}    % repo rate
\newcommand{\Fund}{r^{\mathrm{fun}}}     % unsecured funding rate
\newcommand{\OIS}{r^{\mathrm{ois}}}      % OIS rate
\newcommand{\Coll}{r^{\mathrm{c}}}       % collateral rate
\newcommand{\Rcmt}{R^{\mathrm{cmt}}}     % constant-maturity-treasury rate
\newcommand{\Rswp}{R^{\mathrm{swp}}}     % swap rate (used for risky variant)
\newcommand{\Rec}{\mathrm{Rec}}          % recovery rate
\newcommand{\NPV}{\mathrm{NPV}}          % present-value functional
\newcommand{\PVone}{\mathrm{PV01}}       % risky annuity
\newcommand{\Loss}{\mathrm{Loss}}        % default-leg integral
\newcommand{\Sasw}{S^{\mathrm{ASW}}}     % asset-swap spread
\newcommand{\Scds}{S^{\mathrm{CDS}}}     % CDS spread
\newcommand{\Basis}{\mathrm{Basis}}      % CDS-bond basis

% --- Sensitivities ---
\newcommand{\Diff}[2]{D_{#1}^{#2}}       % Euler's notation: D_x^k[f]
\newcommand{\Riskf}{\mathrm{RiskFactor}} % bond risk factor (= -DV01*1e4)
\newcommand{\DVone}{\mathrm{DV01}}

% --- Indicator / misc ---
\newcommand{\Ind}{\mathbf{1}}
\newcommand{\given}{\,|\,}
\newcommand{\dd}{\,\mathrm{d}}

% --- Greeks-style ---
\newcommand{\Gam}{\Gamma}
\newcommand{\Vega}{\mathcal{V}}


\title{Project Bibliography\\
       \large Source metadata for all papers, verbatim}
\author{}
\date{\today}

\begin{document}
\maketitle

\tableofcontents
\bigskip\hrule\bigskip

% =====================================================================
\section{Pucci (2019) --- Hedging the Treasury Lock}
\label{bib:pucci-2019-tlock}

\subsection*{Identification}
\begin{tabular}{@{}p{3.5cm}p{11cm}@{}}
\toprule
Title           & Hedging the Treasury Lock \\
Author          & Mario Pucci \\
Affiliation     & Banca IMI SpA, Largo Mattioli 3, Milan, 20121, Italy \\
Email           & \texttt{mario.pucci@bancaimi.com} \\
Date            & 19 August 2019 \\
Source          & \href{https://ssrn.com/abstract=3386521}{ssrn:3386521} \\
Local file      & \texttt{papers/pucci\_2019\_tlock/ssrn3386521.pdf} \\
Length          & 25 pages \\
\bottomrule
\end{tabular}

\subsection*{Abstract}
\begin{quote}\itshape
The Treasury lock is a common pre-hedging derivative strategy the
Street offers to their corporate clients. The paper provides a
justification of the common practice of booking a short position in
the Treasury lock as a forward contract on the underlying benchmark
and a short position in the Then-Current Treasury lock as a forward
contract on underlying benchmark rolled over the life of the contract.
\end{quote}

\subsection*{Keywords}
Treasury, lock, T-Lock, hedging, arbitrage-free, convexity, gamma, repo.

\subsection*{Contents}
\begin{enumerate}[noitemsep]
  \item Introduction (Trade Description; Pricing the Standard T-Lock; Approximating a T-Lock as a Forward Contract).
  \item Hedging the Standard T-Lock (Overhedge Error; Delta; Gamma; Replication via Repo).
  \item Trading the Then-Current-T-Lock (Roll P\&L; Forming the New Benchmark).
  \item Statistical Analysis (Overhedge Error Estimation; Negative Gamma Bounds; Assessment of Roll Risk).
  \item Valuation Adjustments.
  \item Conclusions.
  \item Appendices A--D: Simple Math; Key Derivatives; More Math; The Forward Price and the Repo.
\end{enumerate}

\subsection*{References}
22 numbered references. Notable: Hunt \& Kennedy (\emph{Financial Derivatives in Theory and Practice}, 2004); Hull (\emph{Options, Futures and Other Derivatives}, 9th ed., 2015); Taleb (\emph{Dynamic Hedging}, 1996); Apostol (\emph{Calculus}, 1967); Carr \& Madan (\emph{Towards a Theory of Volatility Trading}, 2002); Andersen, Duffie \& Song (\emph{Funding value adjustments}, J. Finance 74:1, 2019); Adams \& Smith (\emph{Pre-Issuance Hedging of Fixed-Rate Debt}, J. Applied Corporate Finance 23:4, 2011); Pucci (\emph{CMT Convexity Correction}, IJTAF 17:8, 2014); Ramirez (\emph{Accounting for Derivatives}, Wiley 2015); Morini \& Prampolini (\emph{Risky Funding}, Landmarks in XVA 2016).

\bigskip\hrule\bigskip

% =====================================================================
\section{Pucci (2014) --- Constant Maturity Treasury Convexity Correction}
\label{bib:pucci-2014-cmt}

\subsection*{Identification}
\begin{tabular}{@{}p{3.5cm}p{11cm}@{}}
\toprule
Title           & Constant Maturity Treasury Convexity Correction \\
Author          & Mario Pucci \\
Affiliation     & Banca IMI SpA, Largo Mattioli 3, Milan, 20121, Italy \\
Email           & \texttt{mario.pucci@bancaimi.com} \\
Dates           & Received 13 Sep 2013; Revised 31 Mar 2014, 4 Jul 2014; Accepted 20 Aug 2014 \\
Published       & \emph{International Journal of Theoretical and Applied Finance} 17(8), 2014 \\
Source          & \href{https://ssrn.com/abstract=3387961}{ssrn:3387961} \\
Local file      & \texttt{papers/pucci\_2014\_cmt/ssrn3387961.pdf} \\
Length          & 15 pages \\
\bottomrule
\end{tabular}

\subsection*{Abstract}
\begin{quote}\itshape
In a Constant Maturity Treasury swap the exotic leg pays, for a given
tenor, the yield-to-maturity computed out of a reference bond curve.
This paper introduces a theoretical framework for the modelling of
constant maturity treasury that takes into account default risk of
bond issuer. As an application, we obtain, under simple but standard
assumptions, analytical convexity corrections for some fundamental
payoffs contingent on the constant maturity treasury.
\end{quote}

\subsection*{Keywords}
Convexity adjustment; default risk; constant maturity.

\subsection*{Contents}
Sections (from document body):
\begin{enumerate}[noitemsep]
  \item Introduction.
  \item Yield To Maturity and Constant Maturity Treasury.
  \item Vanilla payoffs and the forward CMT.
  \item Pricing framework (partial filtration).
  \item Convexity correction under standard assumptions.
  \item A hint on incorporating the smile.
\end{enumerate}

\subsection*{References}
Approximately 12 numbered references. Notable: Hagan (\emph{Convexity Conundrums}, Wilmott 2003); Pelsser (\emph{Mathematical Foundation of Convexity Correction}, Quantitative Finance 2001); Hunt \& Kennedy (Wiley 2004); Boenkost \& Schmidt (notes on convexity and quanto adjustments, 2003); Davies \& Pugachevsky (Risk 2005).

\bigskip\hrule\bigskip

% =====================================================================
\section{Pucci (2012a) --- Constant Maturity Asset Swap Convexity Correction}
\label{bib:pucci-2012-cmasw}

\subsection*{Identification}
\begin{tabular}{@{}p{3.5cm}p{11cm}@{}}
\toprule
Title           & Constant Maturity Asset Swap Convexity Correction \\
Author          & Mario Pucci \\
Affiliation     & Banca IMI, Milan \\
Email           & \texttt{mariopucci@hotmail.it} \\
Date            & 27 February 2012 (revision timestamp) \\
Source          & \href{https://ssrn.com/abstract=1961545}{ssrn:1961545} \\
Local file      & \texttt{papers/pucci\_2012\_cmasw/ssrn1961545.pdf} \\
Length          & 10 pages \\
\bottomrule
\end{tabular}

\subsection*{Abstract}
\begin{quote}\itshape
As a consequence of the high volatility regime recently established
in the government bond market, investors may seek hedging their
exposure to floating asset swap spreads. We obtain an analytical
convexity correction for the asset swap spread which is instrumental
to the pricing of constant maturity asset swaps.
\end{quote}

\subsection*{Keywords}
Not stated explicitly. Topics: asset swap, convexity correction, CMS, constant maturity asset swap.

\subsection*{Contents}
\begin{enumerate}[noitemsep]
  \item Introduction.
  \item Constant Maturity Asset Swaps.
  \item Definitions.
  \item Pricing Framework.
  \item One-Factor Linear Model.
  \item Hedging Considerations.
  \item Numerical Example.
  \item Model Limitations.
  \item Conclusions.
\end{enumerate}

\subsection*{References}
13 numbered references. Notable: Hagan (\emph{Convexity Conundrums}, Wilmott 2003); Pelsser (\emph{Mathematical Foundation of Convexity Correction}, Quantitative Finance 2001); Hunt \& Kennedy (Wiley 2004); Zheng \& Kwok (\emph{Convexity meets replication}, J.\ Futures Markets 31, 2010); Nelson \& Siegel (\emph{Parsimonious modeling of yield curves}, J.\ Business 60, 1987); Lee \& Wang (\emph{Displaced Lognormal Volatility Skews}, Annals of Finance 2009); Choudhry (\emph{CDS basis}, YieldCurve.com 2006); BIS (\emph{Zero-coupon yield curves: technical documentation}, BIS Papers No 25, 2005).

\bigskip\hrule\bigskip

% =====================================================================
\section{Pucci (2012b) --- Pricing Index-Linked Hybrids under a Local Volatility Market Model}
\label{bib:pucci-2012-hybrid}

\subsection*{Identification}
\begin{tabular}{@{}p{3.5cm}p{11cm}@{}}
\toprule
Title           & Pricing Index-Linked Hybrids under a Local Volatility Market Model \\
Author          & Mario Pucci \\
Affiliation     & Banca IMI, Milan \\
Email           & \texttt{mariopucci@hotmail.it} \\
Date            & 22 February 2012 (with May 2012 revision) \\
Source          & \href{https://ssrn.com/abstract=2056277}{ssrn:2056277} \\
Local file      & \texttt{papers/pucci\_2012\_hybrid/ssrn2056277.pdf} \\
Length          & 5 pages \\
\bottomrule
\end{tabular}

\subsection*{Abstract}
\begin{quote}\itshape
We present a methodology to price a cash-settled swaption with strike
contingent on the level of an equity index. The hybrid local
volatility framework allows for calibration to market skew/smile and
straightforward specification of an exogenous correlation. The
framework can be extended to other multi-asset hybrid payoffs of
European type.
\end{quote}

\subsection*{Keywords}
Not stated explicitly. Topics: hybrid product, local volatility,
cash-settled swaption, equity-rates hybrid, index-linked,
correlation.

\subsection*{Contents}
Single-paper format, no explicit TOC. Discusses: hybrid product
pricing problems (unobservable factors, correlation, smile);
proposed local-volatility methodology; calibration to market
skew/smile; extension to other European multi-asset hybrid payoffs.

\subsection*{References}
6 numbered references including: Grzelak \& Oosterlee
(\emph{An Equity-Interest Rate Hybrid Model With Stochastic Volatility
and the Interest Rate Smile}, 2010); Piterbarg
(\emph{Smiling hybrids}, Risk, May 2006).

\bigskip\hrule\bigskip

% =====================================================================
\section{Burgess --- Bond Total Return Swaps: Theory, Pricing \& Practice}
\label{bib:burgess-trs}

\subsection*{Identification}
\begin{tabular}{@{}p{3.5cm}p{11cm}@{}}
\toprule
Title           & Bond Total Return Swaps: Theory, Pricing \& Practice \\
Author          & Nicholas Burgess \\
Affiliation     & Not stated in document \\
Date            & Not stated on cover \\
Source          & \url{https://ssrn.com/author=1728976}; companion repo at \url{https://github.com/nburgessx/SwapsBook} \\
Local file      & \texttt{papers/burgess\_trs/5024091.pdf} \\
Length          & 23 pages \\
\bottomrule
\end{tabular}

\subsection*{Abstract}
No abstract printed.

\subsection*{Keywords}
Not stated. Topics: bond total return swap (TRS); contract specifications; constant units vs.\ constant notional; pricing formulae; worked case study.

\subsection*{Contents}
\textbf{Part One: Theory --- Bond Total Return Swaps.}
\begin{itemize}[noitemsep]
  \item Introduction to Bond TRS.
  \item Why Trade a TRS?
  \item Contract Specifications.
  \item Constant Units vs Constant Notional.
\end{itemize}
\textbf{Part Two: Pricing \& Practice --- Case Studies.}
\begin{itemize}[noitemsep]
  \item Pricing Formulae.
  \item Bond TRS Case Study.
\end{itemize}

\subsection*{References}
No formal numbered reference list is printed. The document links to
the author's SSRN author page (\url{ssrn.com/author=1728976}) and a
companion GitHub repository (\url{github.com/nburgessx/SwapsBook}),
and to the author's book \emph{Low Latency Interest Rate Markets}.

\bigskip\hrule\bigskip

% =====================================================================
\section{Lou (2018) --- Pricing Total Return Swap}
\label{bib:lou-2018-trs}

\subsection*{Identification}
\begin{tabular}{@{}p{3.5cm}p{11cm}@{}}
\toprule
Title           & Pricing Total Return Swap \\
Author          & Wujiang Lou \\
Affiliation     & New York University, Courant Institute of Mathematics \\
Date            & 16 July 2018 \\
Source          & \href{https://ssrn.com/abstract=3217420}{ssrn:3217420} \\
Local file      & \texttt{papers/lou\_2018\_trs/SSRNid32174202.pdf} \\
Length          & 25 pages \\
\bottomrule
\end{tabular}

\subsection*{Abstract}
\begin{quote}\itshape
Total return swap (TRS) involves a pricing dilemma: Libor discounting
of its premium leg forces upfront payment of future funding premium,
and yet replacing Libor with a firm's own funding rate falls into the
well-known FVA debate trap. We consider TRS hedge financing from a
repo market perspective and apply post-crisis derivatives valuation
with collateralization and funding to TRS. We find that the financing
cost of the TRS hedge should be reflected on the security leg, and
the funding premium can only be discounted in conjunction with the
TRS as a whole, depending on margining schemes. An easy to implement,
recursive tree model is developed to value TRS with repo-style
margining or defaultable underlying, together with any value
adjustments.
\end{quote}

\subsection*{Keywords}
Total return swap, TRS pricing dilemma, funding costs, funding
valuation adjustment, liability-side pricing, bond forward.

\subsection*{Contents}
Topics (from inline section headings):
\begin{enumerate}[noitemsep]
  \item Introduction.
  \item TRS mechanics (security leg, funding leg, repo financing).
  \item Pricing under collateralisation and funding.
  \item Recursive tree model for repo-style margining and defaultable underlying.
  \item Value adjustments (XVA components).
  \item Appendix: bond forward derivation under deterministic interest and repo rates.
\end{enumerate}

\subsection*{References}
References to the post-2008 funding / XVA literature including
Andersen, Duffie \& Song (FVA); Burgard \& Kjaer; Piterbarg
(\emph{Funding beyond discounting}). Detailed list verified from
the document body.

\bigskip\hrule\bigskip

% =====================================================================
\section{Zhou (2008) --- Bond Implied CDS Spread and CDS-Bond Basis}
\label{bib:zhou-2008-cds}

\subsection*{Identification}
\begin{tabular}{@{}p{3.5cm}p{11cm}@{}}
\toprule
Title           & Bond Implied CDS Spread and CDS-Bond Basis \\
Author          & Richard Zhou \\
Email           & \texttt{rzhou50@gmail.com} \\
Affiliation     & Not stated; disclaimer notes views are author's own \\
Date            & 15 August 2008 \\
Source          & \href{https://ssrn.com/abstract=1265548}{ssrn:1265548} \\
Local file      & \texttt{papers/zhou\_2008\_cds/ssrn1265548.pdf} \\
Length          & 12 pages \\
\bottomrule
\end{tabular}

\subsection*{Abstract}
\begin{quote}\itshape
We derive a simple formula for calculating the CDS spread implied by
the bond market price. Using no-arbitrage argument, the formula
expresses the bond implied CDS spread as the sum of bond price, bond
coupon and Libor zero curve weighted by risky annuities. We show that
the bond implied CDS spread is consistent with the standard CDS
pricing model if the survival probabilities and recovery are
consistent with the bond price.
\end{quote}

\subsection*{Keywords}
Not stated explicitly. Topics: CDS spread, asset swap spread, CDS-bond basis, survival probability, recovery rate, bond-implied CDS.

\subsection*{Contents}
\begin{enumerate}[noitemsep]
  \item Introduction.
  \item CDS and Asset Swap Spreads (theory; main formulas; model applications).
  \item Numerical examples.
  \item Conclusions.
  \item Appendix A: model derivation.
  \item Appendix B: details.
\end{enumerate}

\subsection*{References}
Numbered list referencing standard credit-derivatives literature including: Davies \& Pugachevsky on bond-implied CDS approximations; Berd et al.\ on survival-based modelling; the standard CDS pricing model literature; references on factors driving CDS-bond basis.

\bigskip\hrule\bigskip

% =====================================================================
\section{Axelsson \& Renström (2013) --- Credit-Linked Notes: Swedish Market Evidence}
\label{bib:axelsson-renstrom-2013-cln}

\subsection*{Identification}
\begin{tabular}{@{}p{3.5cm}p{11cm}@{}}
\toprule
Title           & On the pricing of Credit-Linked Notes: Evidence from the Swedish market \\
Authors         & Marc Axelsson, Tobias Renström \\
Emails          & \texttt{40272@student.hhs.se}, \texttt{40278@student.hhs.se} \\
Affiliation     & Stockholm School of Economics (SSE) --- master's thesis; tutor: Michael Halling, Department of Finance \\
Date            & 15 May 2013 \\
Source          & No SSRN listed on document \\
Local file      & \texttt{papers/axelsson\_renstrom\_2013\_cln/cln\_swedish\_mkt.pdf} \\
Length          & 50 pages \\
\bottomrule
\end{tabular}

\subsection*{Abstract}
\begin{quote}\itshape
Based on international evidence of overpricing, we study the newly
established credit-linked note (CLN) market in Sweden during the
period 2012--2013. We derive a theoretical pricing model based on a
reduced form framework and value the outstanding CLNs in the Swedish
market. Our results suggest significant and consistent overpricing,
ranging from 0.6\% to 6.4\%. Furthermore, when accounting for the
counterparty credit risk exposure in the issuer, the overpricing
increases substantially to a maximum of 8.5\%. Given this considerable
overpricing, it is possible that the counterparty risk is not priced
in the CLNs. Plausibly, investors perceive the issuer as ``too big to
fail'', which is supported by previous findings of Arora et al.\
(2012). Lastly, our study indicates the existence of market frictions
that influence the market pricing mechanism, obfuscating the
relationship between observed prices and credit risk. Based on
findings of Blanco et al.\ (2005), this could be a result of limited
liquidity.
\end{quote}

\subsection*{Keywords}
Structured products, interest certificates, credit-linked notes, reduced form models.

\subsection*{Contents}
Master's-thesis format (introduction, literature review, theory,
data, empirical analysis, conclusion). Major sections expected:
introduction; literature review (Burth et al.\ 2001; Stoimenov \&
Wilkens 2005; Blanco et al.\ 2005; Arora et al.\ 2012); reduced-form
pricing model; data and methodology; empirical results
(overpricing 0.6\%--6.4\% raw, up to 8.5\% with CVA); discussion of
market frictions and liquidity; conclusion.

\subsection*{References}
Extensive thesis-style reference list. Notable: Burth, Kraus \&
Wuhrmann (Swiss structured product pricing, 2001); Stoimenov \&
Wilkens (German equity-linked structured products, 2005); Blanco,
Brennan \& Marsh (CDS-bond basis, 2005); Arora, Gandhi \& Longstaff
(counterparty risk and CDS spreads, 2012); Yiu (2012); Swedberg
(2010).

\bigskip\hrule\bigskip

% =====================================================================
\section{Anonymous --- Treasury Lock Model}
\label{bib:anon-tlock-model}

\subsection*{Identification}
\begin{tabular}{@{}p{3.5cm}p{11cm}@{}}
\toprule
Title           & Treasury Lock Model \\
Author          & Not stated (anonymous model documentation) \\
Affiliation     & Document references \url{https://finpricing.com/lib/FiZeroBond.html}, suggesting FinPricing \\
Date            & Not stated \\
Source          & No SSRN/DOI on document \\
Local file      & \texttt{papers/anon\_treasury\_lock\_model/31670362145195.pdf} \\
Length          & 4 pages \\
\bottomrule
\end{tabular}

\subsection*{Abstract}
No abstract printed. The opening paragraph reads:
\begin{quote}\itshape
Treasury Lock (T-Lock) is a customized agreement that fixes the yield
or price on a specified treasury security for a specific period.
The buyer of the T-Lock is protected from a rise in the yield (i.e.\
a fall in price) of the underlying treasury security during the lock
period. The period or term of the T-Lock is normally from one week to
twelve months.
\end{quote}

\subsection*{Keywords}
Not stated. Topics: Treasury Lock, bond forward pricing, repo curve,
clean/dirty price, locked yield, PV01.

\subsection*{Contents}
\begin{enumerate}[noitemsep]
  \item Bond forward pricing formula (clean price form, repo discount factor).
  \item Model inputs.
  \item T-Lock NPV for clean-price lock.
  \item T-Lock NPV for yield lock (with finite-difference PV01).
  \item T-Lock NPV for dirty-price lock.
  \item Note on PV01/NDELTA as approximate sensitivities.
\end{enumerate}

\subsection*{References}
One external link cited inline: \url{https://finpricing.com/lib/FiZeroBond.html}. No formal reference list.

\bigskip\hrule\bigskip

% =====================================================================
\section{Anonymous --- Why We Don't Discount at the ``Risk-Free Rate''}
\label{bib:anon-discounting}

\subsection*{Identification}
\begin{tabular}{@{}p{3.5cm}p{11cm}@{}}
\toprule
Title           & Why We Don't Discount at the ``Risk-Free Rate'' \\
Subtitle        & What Textbooks Get Wrong About Derivatives Pricing, and What Trading Desks Actually Do \\
Author          & Not stated (anonymous) \\
Date            & Not stated \\
Source          & No SSRN/DOI on document \\
Local file      & \texttt{papers/anon\_discounting\_textbooks/What\_Textbooks\_Get\_Wrong\_About\_Discounting\_1777005506.pdf} \\
Length          & 9 pages \\
\bottomrule
\end{tabular}

\subsection*{Abstract}
No abstract printed. The opening paragraphs read:
\begin{quote}\itshape
If you studied derivatives pricing at university, you were almost
certainly taught that derivatives are priced by taking the expected
payoff under the risk-neutral measure and discounting at the
risk-free rate. \dots\ It's also, in its na\"ive form, wrong. \dots\
The question ``what rate do I discount at?'' turns out to be far
more subtle than the textbooks suggest, and the answer has reshaped
the entire derivatives industry since 2008.
\end{quote}

\subsection*{Keywords}
Not stated. Topics: collateral rate, repo rate, funding rate, CSA,
OIS, multi-curve, XVA (CVA/DVA/FVA), ColVA, partial CSA,
cross-currency basis.

\subsection*{Contents}
\begin{enumerate}[noitemsep]
  \item The General Framework.
  \item The Textbook Argument, and Where It Hides Its Assumptions.
  \item What the Rate Actually Represents: The Repo Rate (incl.\ equity, bond examples; summary table repo vs.\ collateral).
  \item The CSA Revolution.
  \item Collateral Regimes and the Discount Rate (perfect CSA; no CSA; partial CSA).
  \item Linking the Collateral Rate to the Discount Rate (cash flow mechanics; martingale argument; worked swap example; cross-currency CSA; non-cash collateral; decision tree).
  \item From Single-Curve to Multi-Curve.
  \item The Correct Mental Model.
\end{enumerate}

\subsection*{References}
No formal reference list. Hull and Shreve mentioned as textbook examples.

\bigskip\hrule\bigskip

\bigskip\hrule\bigskip

% =====================================================================
\section{Brigo \& Morini (2005) --- CDS Market Formulas and Models}
\label{bib:brigo-morini-2005}

\subsection*{Identification}
\begin{tabular}{@{}p{3.5cm}p{11cm}@{}}
\toprule
Title           & CDS Market Formulas and Models \\
Authors         & Damiano Brigo, Massimo Morini \\
Affiliation     & Brigo: Credit Models, Banca IMI, Milan. Morini: Università di Milano Bicocca \\
Date            & September 2005 \\
Source          & Working paper, no SSRN ID printed on document \\
Local file      & \texttt{papers/brigo\_morini\_2005\_cdsmm/cdsmktfor.pdf} \\
Length          & 28 pages \\
\bottomrule
\end{tabular}

\subsection*{Abstract}
\begin{quote}\itshape
In this work we analyze market payoffs of Credit Default Swaps (CDS)
and we derive rigorous standard market formulas for pricing options
on CDS. Formulas are based on modelling CDS spreads which are
consistent with simple market payoffs, and we introduce a
subfiltration structure allowing all measures to be equivalent to
the risk neutral measure. Then we investigate market CDS spreads
through change of measure and consider possible choices of rates for
modelling a complete term structure of CDS spreads. We also consider
approximations and apply them to pricing of specific market
contracts. Results are derived in a probabilistic framework similar
to that of Jamshidian (2004).
\end{quote}

\subsection*{Keywords}
CDS, CDS option, market model, swap market model, subfiltration,
change of measure, defaultable PV01, Black formula, CMCDS,
constant maturity CDS, conditional independence.

\subsection*{Contents}
\begin{enumerate}[noitemsep]
  \item Introduction.
  \item Credit Default Swaps and Options (par CDS spread; CDS option payoff).
  \item Standard Market Formula for CDS (numeraire $C_{a,b}$, change of measure, dynamics of $R_{a,b}$, empirical application).
  \item Variables for a general CDS Market Model (one-period CDS forward rates; CDS spreads as forward conditional probability ratios; dynamics under different measures).
  \item Dynamics and Constant Maturity CDS (Libor + CDS Market Model under independence; CMCDS pricing with drift-freezing approximation; empirical application on Fiat; alternative model with approximated payoff).
  \item Conclusion.
  \item Appendix (propositions and proofs).
\end{enumerate}

\subsection*{References}
Brigo (2004, 2005), Brigo \& Mercurio (2001), Jamshidian (1997, 2004),
Jeanblanc \& Rutkowski (2000), Hull \& White (2003), Schönbucher
(1999, 2000, 2004), Bielecki \& Rutkowski (2001), Wu (2005), and
the seminal Black (1976) and Brace-Gatarek-Musiela (1997).

\bigskip\hrule\bigskip

% =====================================================================
\section{Ametrano \& Bianchetti (2013) --- Multiple Interest Rate Curve Bootstrapping}
\label{bib:ametrano-bianchetti-2013}

\subsection*{Identification}
\begin{tabular}{@{}p{3.5cm}p{11cm}@{}}
\toprule
Title           & Everything You Always Wanted to Know About Multiple Interest Rate Curve Bootstrapping But Were Afraid To Ask \\
Authors         & Ferdinando M.\ Ametrano, Marco Bianchetti \\
Affiliation     & Ametrano: Financial Engineering, Banca IMI, Milan. Bianchetti: Market Risk Management, Banca Intesa Sanpaolo, Milan \\
Date            & 2 April 2013 (first version 4 February 2013) \\
Source          & \href{https://ssrn.com/abstract=2219548}{ssrn:2219548}.
                  Also a chapter in \emph{Interest Rate Modelling After the Financial Crisis} (Bianchetti \& Morini eds., Risk Books, 2013) \\
Local file      & \texttt{papers/ametrano\_bianchetti\_2013\_multicurve/ssrn2219548.pdf} \\
Length          & 82 pages \\
\bottomrule
\end{tabular}

\subsection*{Abstract}
\begin{quote}\itshape
After the credit and liquidity crisis started in summer 2007 the
market has recognized that multiple yield curves are required for
estimation of both discount and FRA rates with different tenors
(e.g.\ Overnight, Libor 3 months, etc.), consistently with the large
basis spreads and the wide diffusion of bilateral collateral
agreements and central counterparties for derivatives transactions
observed on the market.\dots The theoretical pricing framework is
introduced and modern pricing formulas for plain vanilla interest rate
derivatives, such as Deposits, Forward Rate Agreements (FRA),
Futures, Swaps, OIS, and Basis Swaps, are derived from scratch. The
concrete EUR market case is worked out\dots The implementation of
the proposed algorithms is available open source within the QuantLib
framework.
\end{quote}

\subsection*{Keywords}
crisis, crunch, liquidity risk, funding risk, credit risk, counterparty risk,
collateral, CSA, fixed income, interest rates, negative rates, Libor,
Euribor, Eonia, yield curve, FRA curve, discount curve, single-curve,
multiple-curve, bootstrapping, no-arbitrage, pricing, hedging,
interest rate derivatives, Deposit, FRA, Futures, Swap, IRS, OIS,
Basis Swap, turn of year, spline, greeks, sensitivity, QuantLib.

\subsection*{Contents}
\begin{enumerate}[noitemsep]
  \item Introduction.
  \item Classical and Modern Pricing \& Hedging of Interest Rate Derivatives (single-curve vs.\ multiple-curve approach).
  \item Notation and Basic Theoretical Framework (contract description, funding and collateral, pricing under collateral, interest rates, yield curves).
  \item Bootstrapping Interest Rate Yield Curves (classical vs.\ modern bootstrapping; instruments selection --- Deposits, FRA, Futures, IRS, OIS, Basis Swaps; synthetic instruments; interpolation; negative rates; endogenous vs.\ exogenous; turn of year effect; multiple deltas \& hedging; performance and yield curve checks).
  \item Implementation and Examples of Bootstrapping (Eonia ON, Euribor 1M / 3M / 6M / 12M, basis curves; full EUR market 11-Dec-2012 case).
  \item Conclusions.
  \item Appendices: pricing under collateral, Feynman-Kac, derivations of pricing formulas for FRA, Futures, IRS, OIS, IRBS.
\end{enumerate}

\subsection*{References}
Includes Mercurio (multiple-curve work), Brigo \& Mercurio
(\emph{Interest-Rate Models}), Andersen \& Piterbarg (\emph{Interest
Rate Modeling}), Hagan \& West (curve interpolation), White
(OpenGamma multi-curve), ISDA standard CSA documentation, plus
the open-source QuantLib library which implements the paper's
algorithms.

\bigskip\hrule\bigskip

% =====================================================================
\section{Personal Work --- The Treasury Lock: a comprehensive study}
\label{bib:personal-tlock-comprehensive}

\subsection*{Identification}
\begin{tabular}{@{}p{3.5cm}p{11cm}@{}}
\toprule
Title           & The Treasury Lock --- a comprehensive study \\
Type            & Personal-work integrated study (not a third-party paper) \\
Sources         & (1) Anonymous practitioner T-Lock note; (2) Pucci 2019 (ssrn:3386521) \\
Date assembled  & 2026 \\
Local file      & \texttt{papers/personalWork/tlock\_comprehensive/tlock\_comprehensive.tex} \\
Length          & 10 pages \\
\bottomrule
\end{tabular}

\subsection*{Abstract (own)}
A unified treatment of Treasury Lock pricing combining the operational
pricing recipe of the anonymous practitioner note (three lock
conventions reduced to the bond forward formula) and the theoretical
justification of Pucci~(2019) (payoff identity, forward-contract
overhedge, negative-gamma bound, repo replication, Then-Current
roll). Two parts (Part~I practitioner, Part~II theoretical) plus a
unified validation block reconciling the simple-vs-continuous
compounding conventions on Pucci's 24-Jan-2019 worked example.

\subsection*{Keywords}
T-Lock, repo, forward, convexity, gamma, pre-hedge, overhedge,
Then-Current.

\subsection*{Contents}
\begin{enumerate}[noitemsep]
  \item Introduction and Trade Description.
  \item Unified Notation.
  \item \textbf{Part I --- Practitioner Pricing Model:} bond forward
        building block; three lock conventions; practitioner risks
        and sensitivities.
  \item \textbf{Part II --- Hedging and Convexity Analysis (Pucci 2019):}
        arbitrage-free pricing formula; forward proxy and overhedge;
        greeks (delta, gamma, sign conditions); repo replication;
        Then-Current T-Lock and roll risk; valuation adjustments;
        conclusions.
  \item Validation \& numerical examples (canonical 24-Jan-2019 case).
\end{enumerate}

\subsection*{References}
\begin{itemize}[noitemsep]
\item Anonymous (n.d.), \emph{Treasury Lock Model}, practitioner note.
\item Pucci, M.\ (2019), \emph{Hedging the Treasury Lock}, Banca IMI SpA, SSRN 3386521.
\end{itemize}

\end{document}
